% Generated by analysis/scripts/export_ten_cases.py
\subsection*{统一结果卡}
\begingroup\small
\noindent\begin{tabularx}{\textwidth}{@{}p{24mm}X@{}}\toprule
技术/模式 & MRAM；互补2T2MTJ二状态数字路径；两相写＋双支路终验\\
逻辑粒度/聚合 & B\_S=128 Byte；B\_R=8 Byte。8 Byte对齐组＝64互补bit对；2048事务覆盖矩阵；先两MTJ置P，再择一AP并验两支路\\
资源/相对R0 & 4096活动读对；128方向写支路、64单端绝对P/AP验收lane、128bit状态锁存\\
服务口径 & 无周期维护；终验/返回包括在事务内\\
主导因素/选择 & 一次完整尝试为主点；参考两写槽占60/95 ns；重写对照保留失败尝试时间，非裸脉冲\\
关键对照 & 同外围一次额外整组重写，$\Delta$R=180 ns\\
\bottomrule\end{tabularx}
\vspace{3mm}
\begin{center}\begin{tabular}{@{}lrrr@{}}\toprule
成对条件情景 & $\rho$ (MB/s) & $\tau$ (MB/s) & $\mathrm{RI}^{*}$\\\midrule
短预算 & 99.69 & 142.9 & 0.6978\\
参考（推荐） & 49.81 & 84.21 & 0.5914\\
长预算 & 24.9 & 61.54 & 0.4047\\
\bottomrule\end{tabular}\end{center}
\noindent 参考原始占用：streaming 2570 ns；resident 95 ns。
\par\noindent 范围为有限成对条件情景极值，不是物理不确定性包络。1 MB=$10^6$ Byte；整矩阵16384 Byte=16 KiB。源定位见本例证据笔记；公共基线shared\_baseline。
\endgroup
\clearpage
